\section{Pendahuluan}

% Slide 2: 1.1 Latar Belakang - Limitasi Model Markowitz
\begin{frame}{1.1 Latar Belakang - Limitasi Model Markowitz}
    \begin{itemize}
        \item \textbf{Modern Portfolio Theory:} Model \textit{Mean-Variance} (MV) sebagai fondasi utama alokasi aset dan \textbf{Diversifikasi}.
        \item \textbf{Masalah Utama:}
        \begin{itemize}
            \item Asumsi distribusi \textit{return} yang normal dan stasioner seringkali tidak realistis.
            \item Ketidakmampuan menangkap anomali pasar dan ketergantungan non-linear selama krisis.
        \end{itemize}
        \item \textbf{Solusi:} Transformasi fungsi objektif ke dalam kerangka \textit{Quadratic Unconstrained Binary Optimization} (QUBO).
    \end{itemize}
    \begin{center}
        \textit{Saran Gambar: Grafik Efficient Frontier Klasik vs Realitas Pasar (Fat-tailed distribution).}
    \end{center}
\end{frame}

% Slide 3: 1.1 Latar Belakang - Pendekatan Ekonofisika
\begin{frame}{1.1 Latar Belakang - Pendekatan Ekonofisika}
    \begin{itemize}
        \item \textbf{Model Ising:} Memetakan fungsi biaya portofolio ke dalam struktur Hamiltonian sistem fisik.
        \item \textbf{Representasi Spin:} Aset dianalogikan sebagai \textit{spins} yang berinteraksi dalam kisi kristal.
        \item \textbf{Fenomena Herding:} Sinkronisasi keputusan investor dicerminkan melalui keselarasan orientasi \textit{spin}.
        \item \textbf{Tujuan:} Mencari konfigurasi energi terendah (\textit{ground state}) yang setara dengan alokasi modal paling efisien.
    \end{itemize}
    \begin{center}
        %\includegraphics[width=0.4\textwidth]{Spin_Interaction_each.png} \\
        \textit{Saran Gambar: Visualisasi interaksi spin dalam kisi (Grid/Lattice).}
    \end{center}
\end{frame}

% Slide 4: 1.1 Latar Belakang - Integrasi Teori Permainan
\begin{frame}{1.1 Latar Belakang - Integrasi Teori Permainan}
    \begin{itemize}
        \item \textbf{Exact Potential Game (EPG):} Pemodelan interaksi strategis antar partisipan pasar secara dinamis.
        \item \textbf{Fungsi Potensial ($\Phi$):} Menjamin konsistensi antara utilitas marginal individu dan energi sistem total.
        \item \textbf{Kelebihan:} Deteksi bias aset secara presisi melalui pencarian titik \textit{Nash Equilibrium}.
        \item \textbf{Sinergi:} Integrasi EPG ke model Ising memodifikasi parameter medan lokal ($h_i$) berdasarkan logika permainan.
    \end{itemize}
\end{frame}

% Slide 5: 1.1 Latar Belakang - Strategi Warm-Start
\begin{frame}{1.1 Latar Belakang - Strategi Warm-Start}
    \begin{itemize}
        \item \textbf{Barren Plateaus:} Kendala lenyapnya gradien pada sirkuit kuantum variasional yang menghambat optimasi.
        \item \textbf{Warm-Start Strategis:} Memanfaatkan konfigurasi \textit{Nash Equilibrium} sebagai inisialisasi parameter awal VQE.
        \item \textbf{Mekanisme:} Menempatkan sirkuit pada wilayah probabilitas tinggi sejak iterasi pertama.
        \item \textbf{Hasil:} Mempercepat konvergensi energi menuju solusi optimal secara signifikan.
    \end{itemize}
\end{frame}

% Slide 6: 1.1 Latar Belakang - VQE dan Optimasi SPSA
\begin{frame}{1.1 Latar Belakang - VQE dan Optimasi SPSA}
    \begin{itemize}
        \item \textbf{Variational Quantum Eigensolver (VQE):} Algoritma hibrida yang fleksibel untuk perangkat keras era NISQ.
        \item \textbf{Optimasi SPSA:}
        \begin{itemize}
            \item Estimasi gradien efisien (hanya 2 pengukuran per iterasi).
            \item Tangguh terhadap \textit{noise} pengukuran dan fluktuasi statistik pasar.
        \end{itemize}
        \item \textbf{Ansatz EfficientSU(2):} Struktur sirkuit ringkas yang meminimalkan akumulasi \textit{gate noise} tanpa mengurangi ekspresivitas.
    \end{itemize}
\end{frame}

% Slide 7: 1.2 Rumusan Masalah
\begin{frame}{1.2 Rumusan Masalah}
    \begin{enumerate}
        \item Bagaimana perumusan bias aset ($h_i$) dan kopling ($J_{ij}$) menggunakan fungsi potensial EPG memengaruhi akurasi pemetaan Hamiltonian?
        \item Bagaimana efektivitas \textit{Nash Equilibrium} sebagai strategi \textit{warm-start} dalam mengatasi kendala inisialisasi algoritma VQE?
        \item Bagaimana performa algoritma VQE dengan optimasi SPSA dan \textit{Ansatz} EfficientSU(2) dalam menemukan solusi \textit{ground state}?
    \end{enumerate}
\end{frame}

% Slide 8: 1.3 Tujuan Penelitian
\begin{frame}{1.3 Tujuan Penelitian}
    \begin{enumerate}
        \item Menganalisis formulasi bias dan interaksi aset melalui pendekatan fungsi potensial EPG dalam sistem Hamiltonian.
        \item Mengevaluasi kontribusi strategi \textit{warm-start} berbasis \textit{Nash Equilibrium} dalam meningkatkan efisiensi pencarian solusi VQE.
        \item Menguji kapabilitas kombinasi VQE, SPSA, dan EfficientSU(2) untuk alokasi aset di pasar yang kompleks.
    \end{enumerate}
\end{frame}

% Slide 9: 1.4 Batasan Masalah
\begin{frame}{1.4 Batasan Masalah}
    \begin{itemize}
        \item \textbf{Portofolio:} Pemilihan aset dari pasangan sektor berbeda ($N=2$ dan $N=4$).
        \item \textbf{Model Game:} Skema \textit{non zero-sum game} dalam kerangka \textit{Exact Potential Game}.
        \item \textbf{Algoritma:} Pencarian ekuilibrium menggunakan \textit{Sequential Best Response} (SBR).
        \item \textbf{Implementasi:} Simulasi menggunakan framework \textit{PennyLane}.
    \end{itemize}
\end{frame}

% Slide 10: 1.5 Manfaat Penelitian
\begin{frame}{1.5 Manfaat Penelitian}
    \begin{itemize}
        \item Memberikan alternatif metode optimasi portofolio yang lebih tangguh terhadap dinamika korelasi pasar.
        \item Berkontribusi dalam literatur akademis mengenai penerapan komputasi kuantum pada bidang finansial (Econophysics).
        \item Menjembatani logika teori permainan dengan presisi optimasi mekanika kuantum.
    \end{itemize}
\end{frame}
